\chapter{Moles}

\section*{Definition and Overview}
Moles, clinically referred to as melanocytic naevi (singular: naevus), are common, benign skin lesions composed of nests of proliferating melanocytes (the specialised, pigment-producing cells of the epidermis). Moles can be congenital (present at birth or developing within the first year of life) or acquired (appearing during childhood, adolescence, and young adulthood). While the vast majority of melanocytic naevi are completely benign and remain stable throughout a person's lifetime, they are clinically significant because they can mimic malignant melanoma, and in rare instances, a pre-existing mole can undergo malignant transformation into melanoma.

\section*{Epidemiology}
Melanocytic naevi are highly prevalent, particularly in light-skinned populations. In Australia, which has one of the highest incidences of skin cancer globally, the average adult has between 10 and 40 moles. The number of acquired moles increases during childhood, peaks in the third decade of life, and gradually declines or undergoes regression (involution) in older age. The development of a high number of moles, as well as atypical or dysplastic naevi, is strongly associated with genetic factors (such as a family history of atypical mole syndrome), fair skin type, light-coloured eyes, red or blonde hair, and chronic or intermittent intense ultraviolet (UV) radiation exposure, particularly during childhood.

\section*{Aetiology and Pathophysiology}
The formation of moles is driven by genetic mutations and environmental stimulation of melanocytes:
\begin{itemize}
    \item \textbf{Clonal melanocyte proliferation:} Moles occur when melanocytes proliferate in clusters or "nests" along the dermo-epidermal junction (junctional naevi), within the dermis (intradermal naevi), or in both areas (compound naevi).
    \item \textbf{Somatic BRAF and NRAS mutations:} Activating mutations in the proto-oncogene \textit{BRAF} (most commonly the BRAF V600E mutation) or the \textit{NRAS} gene are present in up to 80\% of benign acquired moles, driving initial proliferation followed by oncogene-induced senescence, which halts further growth.
    \item \textbf{Ultraviolet (UV) radiation:} Exposure to UV radiation from the sun damages cellular DNA, promoting melanocytic proliferation and increasing the risk of somatic mutations that lead to atypical naevi or melanoma.
    \item \textbf{Hormonal influences:} Melanocytic naevi are responsive to hormonal changes; they frequently darken, enlarge, or become more numerous during puberty, pregnancy, or systemic corticosteroid therapy.
\end{itemize}

\section*{Clinical Features}
Benign moles typically exhibit symmetry, uniform colouring, and stable borders, whereas atypical moles require careful monitoring:
\begin{itemize}
    \item \textbf{Common acquired naevi:} Typically small (less than 6 millimetres in diameter), round or oval, symmetrical, with well-demarcated borders and a uniform colour (ranging from pink or flesh-coloured to tan, brown, or black). They may be flat (macular) or raised (papular or nodular).
    \item \textbf{Dysplastic (atypical) naevi:} Usually larger (greater than 6 millimetres), with asymmetrical shapes, irregular or smudged borders, and variable pigmentation (a mix of pink, tan, brown, and black).
    \item \textbf{The ABCDE checklist:} Used to identify suspicious lesions: \textbf{A}symmetry, \textbf{B}order irregularity, \textbf{C}olour variation (multiple shades), \textbf{D}iameter greater than 6 millimetres, and \textbf{E}volving (any change in size, shape, colour, or new symptoms like itching, crusting, or bleeding).
    \item \textbf{The "ugly duckling" sign:} A clinical concept where a mole looks significantly different from all other moles on an individual's body, indicating a higher risk of malignancy.
\end{itemize}

\section*{Differential Diagnosis}
Suspicious moles must be distinguished from other pigmented and non-pigmented skin lesions:
\begin{itemize}
    \item \textbf{Malignant melanoma:} A primary skin cancer that can arise de novo (in normal skin) or, less commonly, from a pre-existing mole, presenting with irregular, asymmetric, and changing features.
    \item \textbf{Seborrheic keratosis:} A benign, non-melanocytic epidermal lesion with a characteristic "stuck-on" appearance, waxy or scaly surface, and multiple comedo-like openings, common in older adults.
    \item \textbf{Dermatofibroma:} A benign dermal nodule, most common on the lower limbs, characterised by a positive "dimple sign" (puckering inward when the surrounding skin is pinched).
    \item \textbf{Pigmented basal cell carcinoma:} A common, slow-growing skin cancer that can mimic a dark mole but typically exhibits a pearly translucent border, telangiectasia (fine blood vessels), and central ulceration.
    \item \textbf{Solar lentigines:} Flat, brown macules (commonly called age spots or liver spots) that develop on sun-exposed areas but lack the histological nests of melanocytes characteristic of true naevi.
\end{itemize}

\section*{When to Seek Medical Care}
Individuals should seek a professional skin check if they notice any new moles developing in adulthood, or if an existing mole changes in size, shape, or colour, or begins to bleed, itch, or ulcerate.

\textbf{Call 000 (or local emergency services) for:}
\begin{itemize}
    \item \textbf{Severe bleeding:} Rapid, uncontrolled bleeding from a skin lesion or biopsy site.
    \item \textbf{Anaphylaxis:} A severe, systemic allergic reaction (difficulty breathing, swelling of the tongue/throat, dizziness) following the injection of local anaesthetics during a mole removal.
    \item \textbf{Spreading local infection:} High fever, chills, rapid spreading redness, severe pain, and pus draining from a wound after a mole biopsy or excision.
\end{itemize}

\section*{Diagnosis and Investigations}
Diagnosis is established through clinical skin checks, specialised imaging, and histopathological analysis:
\begin{itemize}
    \item \textbf{Dermoscopy:} Non-invasive evaluation of pigmented lesions using a handheld dermatoscope, allowing visualisation of sub-epidermal structures and pigment networks.
    \textbf{Full-body skin examination:} A thorough check of all skin surfaces—including the scalp, ears, buttocks, soles of the feet, and between the toes—by a trained medical practitioner.
    \item \textbf{Excisional biopsy:} The gold standard for diagnosing a suspicious mole; the entire lesion is removed with a 2-millimetre margin of normal skin and sent for histopathological examination. Incisional or punch biopsies are generally discouraged for suspected melanoma.
    \item \textbf{Sequential digital dermoscopy imaging (SDDI):} Serial high-resolution photography of moles in high-risk patients to detect micro-changes over a 3- to 12-month period.
\end{itemize}

\section*{Management}
Stable, benign moles require no treatment, whereas suspicious or symptomatic moles must be managed surgically:
\begin{itemize}
    \item \textbf{Conservative monitoring:} Regular self-examination and professional skin checks for stable, benign naevi.
    \item \textbf{Surgical excision:} Removal under local anaesthesia using an elliptical excision, indicated for moles showing signs of atypia, those constantly irritated by clothing, or for cosmetic reasons.
    \item \textbf{Histopathological analysis:} All excised naevi must be examined microscopically to confirm the diagnosis and ensure clear margins if atypical or malignant.
    \item \textbf{Avoidance of destructive methods:} Cryotherapy, laser ablation, and electrocautery should not be used to remove melanocytic moles, as they destroy the tissue, preventing histopathological analysis and risking delayed diagnosis of melanoma.
\end{itemize}

\section*{Complications}
The primary complications associated with moles relate to malignancy and surgical procedures:
\begin{itemize}
    \item \textbf{Malignant transformation:} The progression of a pre-existing naevus (particularly a congenital or dysplastic naevus) into invasive melanoma, which can metastasise via lymphatic and haematogenous pathways.
    \item \textbf{Surgical complications:} Bleeding, local wound infection, dehiscence (opening of the wound), and scarring (including hypertrophic or keloid scar formation).
    \item \textbf{Psychological anxiety:} Significant distress and anxiety in patients with numerous or atypical moles regarding their skin cancer risk.
\end{itemize}

\section*{Prognosis and Outcomes}
The prognosis for individuals with benign melanocytic naevi is excellent, as these lesions do not cause physical harm or reduce life expectancy. For atypical moles, the vast majority will never progress to melanoma. In cases where a mole does transform into melanoma, the prognosis is highly dependent on early detection. Melanomas diagnosed at an early stage (Breslow thickness less than 1 millimetre) have a 5-year survival rate exceeding 98\% in Australia. However, advanced or metastatic melanoma carries a significantly poorer prognosis, emphasizing the importance of surveillance.

\section*{Prevention and Self-Care}
Preventative and self-care measures are critical to reducing the risk of mole changes and skin cancer:
\begin{itemize}
    \item \textbf{Sun-safe behaviours (the five S's):} Slip on sun-protective clothing, Slop on SPF 30+ or higher broad-spectrum sunscreen, Slap on a broad-brimmed hat, Seek shade, and Slide on wrap-around sunglasses.
    \item \textbf{Monthly skin self-examinations:} Visually checking the entire body once a month using mirrors, paying close attention to hard-to-see areas like the back and scalp.
    \item \textbf{Avoid artificial UV exposure:} Never using solariums or tanning beds, which deliver concentrated UV radiation and significantly increase melanoma risk.
    \item \textbf{Professional monitoring:} Scheduling annual skin checks with a GP or dermatologist, especially for individuals with more than 100 moles, atypical moles, or a personal/family history of skin cancer.
\end{itemize}

\section*{Sources and Further Reading}
The following organisations provide reliable, evidence-based guidelines and resources on moles and skin cancer prevention:
\begin{itemize}
    \item Cancer Council Australia. \textit{Check for signs of skin cancer.} https://www.cancer.org.au/cancer-information/causes-and-prevention/sun-safety/check-for-signs-of-skin-cancer
    \item Healthdirect Australia. \textit{Moles.} https://www.healthdirect.gov.au/moles
    \item NHS. \textit{Moles.} https://www.nhs.uk/conditions/moles/
    \item Mayo Clinic. \textit{Moles - Symptoms and causes.} https://www.mayoclinic.org/diseases-conditions/moles/symptoms-causes/syc-20375200
    \item Australasian College of Dermatologists. \textit{Moles and Melanoma.} https://www.dermcoll.edu.au/atoz/moles-and-melanoma/
    \item DermNet NZ. \textit{Melanocytic naevi.} https://dermnetnz.org/topics/melanocytic-naevus
\end{itemize}
